% GENERATED FILE. Edit canonical lessons, then run npm run generate:books.
% canonical-source-hash: fnv1a64:3c0ec9468ee3cb88
% canonical-lessons: GU-C547-albatt, GU-C547-hargij, GU-C547-chata, GU-C547-balke, GU-C547-viprit, GU-R547-review-541-547

\chapter{Even So}
\label{ch:gu-even-so}
\hlchaptermodality{547}

\begin{chapteropening}
\textbf{By the end of this chapter:} \emph{I can add a contrast to a condition with despite, rather, of course and by no means.}

Complete the last lesson of chapter 547: I can add a contrast to a condition with despite, rather, of course and by no means.
\end{chapteropening}

\section[\texorpdfstring{\gu{અલબત્ત} (albatt)}{અલબત્ત (albatt)}]{\gu{અલબત્ત} (albatt) — of course (granted)}

\label{lesson:GU-C547-albatt}



\begin{quote}
Before the new one: say the Gujarati for in addition to, then the Gujarati for ever.
\end{quote}

\subsection*{You'll want to know: \gu{અલબત્ત}}
\textbf{\gu{અલબત્ત}} — \emph{albatt} — \textquotedblleft{}of course (granted)\textquotedblright{}.

\begin{quote}
\gu{જો} \gu{તમે} \gu{બોલાવશો} \gu{તો} \gu{અલબત્ત} \gu{હું} \gu{આવીશ}.
\end{quote}

\begin{quote}
\emph{jo tame bolāvśo to albatt hũ āvīś.}
\end{quote}

\begin{quote}
If you invite me, of course I'll come.
\end{quote}

\begin{cousinweb}
\textbf{\gu{અલબત્ત}} (\emph{albatt}) is usually traced to Persian \emph{albatta}, \textquotedblleft{}certainly\textquotedblright{}, a word Persian took from Arabic.
\end{cousinweb}

\begin{grammarlens}[title={What you've built}]
One more word for saying what happens if.
\end{grammarlens}

\subsection*{Your turn}
\begin{itemize}
\raggedright
  \item \emph{Say it:} \emph{albatt}
  \item \emph{Say it:} \emph{albatt}, once more
  \item \emph{Say it:} \emph{jo tame bolāvśo to albatt hũ āvīś.} — If you invite me, of course I'll come.
  \item \emph{Recall:} say the Gujarati for in addition to, then the Gujarati for ever, then say \emph{albatt} again
\end{itemize}

\begin{tcolorbox}[breakable,colback=teal!4,colframe=teal!35!black,title={Before you move on}]
What does \gu{અલબત્ત} mean? (\textquotedblleft{}Of course (granted)\textquotedblright{}.) Say it once more.
\end{tcolorbox}
\section[\texorpdfstring{\gu{હરગિજ} (hargij)}{હરગિજ (hargij)}]{\gu{હરગિજ} (hargij) — by no means (with a negative)}

\label{lesson:GU-C547-hargij}



\begin{quote}
Before the new one: say the Gujarati for ever, then the Gujarati for of course.
\end{quote}

\subsection*{You'll want to know: \gu{હરગિજ}}
\textbf{\gu{હરગિજ}} — \emph{hargij} — \textquotedblleft{}by no means (with a negative)\textquotedblright{}.

\begin{quote}
\gu{જો} \gu{વરસાદ} \gu{પડશે} \gu{તો} \gu{હું} \gu{હરગિજ} \gu{બહાર} \gu{નહીં} \gu{જાઉં}.
\end{quote}

\begin{quote}
\emph{jo varsād paḍśe to hũ hargij bahār nahī̃ jāũ.}
\end{quote}

\begin{quote}
If it rains, I will by no means go out.
\end{quote}

\begin{cousinweb}
\textbf{\gu{હરગિજ}} (\emph{hargij}) is usually traced to Persian \emph{hargiz}, \textquotedblleft{}ever\textquotedblright{}, which with a negative means \textquotedblleft{}never\textquotedblright{}; in Gujarati it comes with \textbf{\gu{નહીં}}.
\end{cousinweb}

\begin{grammarlens}[title={What you've built}]
One more word for saying what happens if.
\end{grammarlens}

\subsection*{Your turn}
\begin{itemize}
\raggedright
  \item \emph{Say it:} \emph{hargij}
  \item \emph{Say it:} \emph{hargij}, once more
  \item \emph{Say it:} \emph{jo varsād paḍśe to hũ hargij bahār nahī̃ jāũ.} — If it rains, I will by no means go out.
  \item \emph{Recall:} say the Gujarati for ever, then the Gujarati for of course, then say \emph{hargij} again
\end{itemize}

\begin{tcolorbox}[breakable,colback=teal!4,colframe=teal!35!black,title={Before you move on}]
What does \gu{હરગિજ} mean? (\textquotedblleft{}By no means (with a negative)\textquotedblright{}.) Say it once more.
\end{tcolorbox}
\section[\texorpdfstring{\gu{છતાં} (chatā̃)}{છતાં (chatā̃)}]{\gu{છતાં} (chatā̃) — despite (after a noun)}

\label{lesson:GU-C547-chata}



\begin{quote}
Before the new one: say the Gujarati for of course, then the Gujarati for by no means.
\end{quote}

\subsection*{You'll want to know: \gu{છતાં}}
\textbf{\gu{છતાં}} — \emph{chatā̃} — \textquotedblleft{}despite (after a noun)\textquotedblright{}.

\begin{quote}
\gu{જો} \gu{તાવ} \gu{છતાં} \gu{તે} \gu{શાળાએ} \gu{જશે} \gu{તો} \gu{તે} \gu{વધારે} \gu{બીમાર} \gu{થશે}.
\end{quote}

\begin{quote}
\emph{jo tāv chatā̃ te śāḷāe jaśe to te vadhāre bīmār thaśe.}
\end{quote}

\begin{quote}
If he goes to school despite the fever, he'll get worse.
\end{quote}

\begin{cousinweb}
\textbf{\gu{છતાં}} (\emph{chatā̃}) follows its noun: \textbf{\gu{તાવ} \gu{છતાં}}, despite the fever. No dictionary we use records an older form, so its history is not settled.
\end{cousinweb}

\begin{grammarlens}[title={What you've built}]
One more word for saying what happens if.
\end{grammarlens}

\subsection*{Your turn}
\begin{itemize}
\raggedright
  \item \emph{Say it:} \emph{chatā̃}
  \item \emph{Say it:} \emph{chatā̃}, once more
  \item \emph{Say it:} \emph{jo tāv chatā̃ te śāḷāe jaśe to te vadhāre bīmār thaśe.} — If he goes to school despite the fever, he'll get worse.
  \item \emph{Recall:} say the Gujarati for of course, then the Gujarati for by no means, then say \emph{chatā̃} again
\end{itemize}

\begin{tcolorbox}[breakable,colback=teal!4,colframe=teal!35!black,title={Before you move on}]
What does \gu{છતાં} mean? (\textquotedblleft{}Despite (after a noun)\textquotedblright{}.) Say it once more.
\end{tcolorbox}
\section[\texorpdfstring{\gu{બલકે} (balke)}{બલકે (balke)}]{\gu{બલકે} (balke) — rather (on the contrary)}

\label{lesson:GU-C547-balke}



\begin{quote}
Before the new one: say the Gujarati for by no means, then the Gujarati for despite.
\end{quote}

\subsection*{You'll want to know: \gu{બલકે}}
\textbf{\gu{બલકે}} — \emph{balke} — \textquotedblleft{}rather (on the contrary)\textquotedblright{}.

\begin{quote}
\gu{જો} \gu{તમે} \gu{આરામ} \gu{કરશો} \gu{તો} \gu{કામ} \gu{નહીં} \gu{બગડે}, \gu{બલકે} \gu{વધારે} \gu{સારું} \gu{થશે}.
\end{quote}

\begin{quote}
\emph{jo tame ārām karśo to kām nahī̃ bagḍe, balke vadhāre sārũ thaśe.}
\end{quote}

\begin{quote}
If you rest, the work won't suffer; rather, it will go better.
\end{quote}

\begin{cousinweb}
\textbf{\gu{બલકે}} (\emph{balke}) is usually traced to Persian \emph{balki}, \textquotedblleft{}but rather\textquotedblright{}.
\end{cousinweb}

\begin{grammarlens}[title={What you've built}]
One more word for saying what happens if.
\end{grammarlens}

\subsection*{Your turn}
\begin{itemize}
\raggedright
  \item \emph{Say it:} \emph{balke}
  \item \emph{Say it:} \emph{balke}, once more
  \item \emph{Say it:} \emph{jo tame ārām karśo to kām nahī̃ bagḍe, balke vadhāre sārũ thaśe.} — If you rest, the work won't suffer; rather, it will go better.
  \item \emph{Recall:} say the Gujarati for by no means, then the Gujarati for despite, then say \emph{balke} again
\end{itemize}

\begin{tcolorbox}[breakable,colback=teal!4,colframe=teal!35!black,title={Before you move on}]
What does \gu{બલકે} mean? (\textquotedblleft{}Rather (on the contrary)\textquotedblright{}.) Say it once more.
\end{tcolorbox}
\section[\texorpdfstring{\gu{વિપરીત} (viprīt)}{વિપરીત (viprīt)}]{\gu{વિપરીત} (viprīt) — adverse (going against you)}

\label{lesson:GU-C547-viprit}



\begin{quote}
Before the new one: say the Gujarati for despite, then the Gujarati for rather.
\end{quote}

\subsection*{You'll want to know: \gu{વિપરીત}}
\textbf{\gu{વિપરીત}} — \emph{viprīt} — \textquotedblleft{}adverse (going against you)\textquotedblright{}.

\begin{quote}
\gu{જો} \gu{હવામાન} \gu{વિપરીત} \gu{હશે} \gu{તો} \gu{અમે} \gu{પ્રવાસ} \gu{ઠેલીશું}.
\end{quote}

\begin{quote}
\emph{jo havāmān viprīt haśe to ame pravās ṭhelīśũ.}
\end{quote}

\begin{quote}
If the weather turns against us, we will put off the trip.
\end{quote}

\begin{cousinweb}
\textbf{\gu{વિપરીત}} (\emph{viprīt}) is a tatsama, taken from Sanskrit \emph{viparīta}, \textquotedblleft{}turned round; contrary\textquotedblright{} (Monier-Williams). Said of weather or of luck, it means going against you.
\end{cousinweb}

\begin{grammarlens}[title={What you've built}]
One more word for saying what happens if.
\end{grammarlens}

\subsection*{Your turn}
\begin{itemize}
\raggedright
  \item \emph{Say it:} \emph{viprīt}
  \item \emph{Say it:} \emph{viprīt}, once more
  \item \emph{Say it:} \emph{jo havāmān viprīt haśe to ame pravās ṭhelīśũ.} — If the weather turns against us, we will put off the trip.
  \item \emph{Recall:} say the Gujarati for despite, then the Gujarati for rather, then say \emph{viprīt} again
\end{itemize}

\begin{tcolorbox}[breakable,colback=teal!4,colframe=teal!35!black,title={Before you move on}]
What does \gu{વિપરીત} mean? (\textquotedblleft{}Adverse (going against you)\textquotedblright{}.) Say it once more.
\end{tcolorbox}
\section[(dialogue)]{Looking back: from Promises and Trust to Even So}

\label{lesson:GU-R547-review-541-547}



\begin{quote}
Say the Gujarati for a solemn oath and an undertaking.
\end{quote}

\subsection*{Your turn}
Say these aloud:

\begin{itemize}
\raggedright
  \item a solemn oath, an undertaking, confidence, a loophole, a proposal
  \item a compromise, negotiations, a mediator, a recommendation, a polite request
  \item an exception, suitable, deserving, capable, special
  \item a difference, a distinction, separate, a copy, evenly matched
  \item effortless, complicated, importance, a try-out, a clever trick
  \item without, without (formal), except, in addition to, ever
  \item of course, by no means, despite, rather, adverse
\end{itemize}

\begin{tcolorbox}[breakable,colback=teal!4,colframe=teal!35!black,title={Before you move on}]
A solemn oath? (\textbf{\gu{શપથ}}, \emph{śapath}.) An undertaking? (\textbf{\gu{વાયદો}}, \emph{vāydo}.) Confidence? (\textbf{\gu{ભરોસો}}, \emph{bharoso}.) A loophole? (\textbf{\gu{છટકબારી}}, \emph{chaṭakbārī}.) A proposal? (\textbf{\gu{પ્રસ્તાવ}}, \emph{prastāv}.) A compromise? (\textbf{\gu{સમાધાન}}, \emph{samādhān}.) Negotiations? (\textbf{\gu{વાટાઘાટ}}, \emph{vāṭāghāṭ}.) A mediator? (\textbf{\gu{મધ્યસ્થી}}, \emph{madhyasthī}.) A recommendation? (\textbf{\gu{ભલામણ}}, \emph{bhalāmaṇ}.) A polite request? (\textbf{\gu{વિનંતી}}, \emph{vinantī}.) An exception? (\textbf{\gu{અપવાદ}}, \emph{apvād}.) Suitable? (\textbf{\gu{યોગ્ય}}, \emph{yogya}.) Deserving? (\textbf{\gu{લાયક}}, \emph{lāyak}.) Capable? (\textbf{\gu{કાબેલ}}, \emph{kābel}.) Special? (\textbf{\gu{ખાસ}}, \emph{khās}.) A difference? (\textbf{\gu{ફેર}}, \emph{pher}.) A distinction? (\textbf{\gu{તફાવત}}, \emph{taphāvat}.) Separate? (\textbf{\gu{અલગ}}, \emph{alag}.) A copy? (\textbf{\gu{નકલ}}, \emph{nakal}.) Evenly matched? (\textbf{\gu{બરાબર}}, \emph{barābar}.) Effortless? (\textbf{\gu{આસાન}}, \emph{āsān}.) Complicated? (\textbf{\gu{જટિલ}}, \emph{jaṭil}.) Importance? (\textbf{\gu{મહત્ત્વ}}, \emph{mahattva}.) A try-out? (\textbf{\gu{અજમાયશ}}, \emph{ajmāyaś}.) A clever trick? (\textbf{\gu{યુક્તિ}}, \emph{yukti}.) Without? (\textbf{\gu{વગર}}, \emph{vagar}.) Without (formal)? (\textbf{\gu{વિના}}, \emph{vinā}.) Except? (\textbf{\gu{સિવાય}}, \emph{sivāy}.) In addition to? (\textbf{\gu{ઉપરાંત}}, \emph{uprā̃t}.) Ever? (\textbf{\gu{કદી}}, \emph{kadī}.) Of course? (\textbf{\gu{અલબત્ત}}, \emph{albatt}.) By no means? (\textbf{\gu{હરગિજ}}, \emph{hargij}.) Despite? (\textbf{\gu{છતાં}}, \emph{chatā̃}.) Rather? (\textbf{\gu{બલકે}}, \emph{balke}.) Adverse? (\textbf{\gu{વિપરીત}}, \emph{viprīt}.)
\end{tcolorbox}
